\begin{abstract}
Fleet managers for warehouse mobile robots must allocate a stream of
pickup-and-delivery tasks, plan collision-free paths for every robot and keep
batteries from running out, all in a loop that never stops. Many fleet
managers connect these problems through free-space distance estimates. Prior
work couples allocation and path finding exactly for one-shot instances, and
approximately in online pickup and delivery by pricing tasks with planned
paths or with the delay that existing reservations impose. In this paper we
introduce \emph{\sname}, an end-to-end lifelong fleet manager for grid
warehouses that uses a cheap member of that family: assignment costs are
arrival times at the pickup of the same windowed space-time A* that plans the
fleet, evaluated against the reservations of the robots that are already
busy, and fed to a Hungarian assignment. A hold cascade, a variant of
existing fail policies that re-plans the robots whose reservations cross a
held cell, keeps prioritized planning conflict-free by construction, and
token-passing endpoint capacities bound dock queues. We prove that
every executed step is vertex- and swap-conflict-free by construction, and
restate the standard optimality of the Hungarian step for its cost matrix;
progress is not guaranteed, and battery
safety holds only under delay bounds that we state but cannot enforce, so we
measure both. We evaluate the whole stack in a seeded discrete-time
simulator with Poisson task streams, loading dwell and an explicit battery
model on four layouts, one of them generated from the zone and dock labels of
an operator-console wireframe and one with one-lane aisles. An independent trajectory
auditor finds no executed conflict in 4{,}740 runs with 60.4 million
robot-steps; with normal batteries no robot depletes and every run delivers
all of its tasks, but with doubled drain the charging policy strands robots
in two thirds of the runs. Without a repair of the reservations that a hold
invalidates, conflicts are executed, and without dynamic priorities the fleet gridlocks. The controlled study of
this pickup-leg coupling, in a loop where dock capacity limits the tasks in
flight, is a negative result that we report as such: greedy, Hungarian,
congestion-proxy and reservation-aware pickup costs stay within 2.9\% of each other
in mean throughput in every one of twelve layout/fleet/load cells, and four
regimes chosen to favour the coupling and committed before their results
(one-lane aisles, hotspot demand, long empty trips, saturated chargers) show
no resolved gain of the reservation-aware cost over greedy assignment in any
regime, and over Hungarian assignment only in the regime in which every method
strands robots; the proxy shows only small effects of either sign outside that
regime.
The study detects no large effect of the cost model but is not an
equivalence test, and it does not cover costs that price the whole trip or
the dock queue. Token passing and token
passing with task swaps deliver 16--66\% less than our loop at its default
dock capacity, mainly because their endpoint rule allows one robot in
flight per dock: with that same limit, token passing with task swaps comes
within 2.5\% of our loop and token passing stays 11--26\% behind. Dock
capacity ($+26$--$29\%$), the planning window
($\pm7\%$) and an optional windowed CBS backend (up to $-8\%$ makespan on small
fleets) are the levers that move the fleet. Code, results and this paper are at
\url{https://github.com/kjungmo/robo-ops}.
\end{abstract}
